\documentclass[DIN, pagenumber=false, parskip=half,
               fromalign=right, fromphone=false, fromfax=false, frombank=true,
               yourmail=false,
               yourref=false,
               fromrule=false]{scrlttr2}


%\usepackage{ngerman}
\usepackage[utf8]{inputenc}
\usepackage[T1]{fontenc}
\usepackage{textcomp}
\usepackage{longtable}
\usepackage{eurosym}

\LoadLetterOption{template}

\newcommand{\stdlohn}[0]{11 \euro}

\renewcommand*{\raggedsignature}{\raggedright}

\setkomavar{fromname}{GNUnet e. V.}
\setkomavar{fromaddress}{Boltzmannstraße 3 \\ 85748 Garching bei München}
\setkomavar{frombank}{IBAN: DE67830654080004822650 \\ BIC: GENODEF1SLR}

\setkomavar{subject}[]{1JD776.003884}
%\setkomavar{yourmail}[Ihre Angebots-Nr.]{2011012631}
%\setkomavar{yourref}[Ihre Kundennummer]{263}
\setkomavar{date}[Datum]{\today}


%===================================
% footer
\firstfoot{
	\parbox[t]{\textwidth}{\footnotesize

	\begin{tabular}[t]{l@{}}
		 \multicolumn{1}{@{}l@{}}{Association:}\\
		 GNUnet e. V. \\
		 Boltzmannstraße 3, 85748 Garching bei München \\ \\
		 E-Mail: ev@gnunet.org \\
	\end{tabular}
	\hfill{}
	\begin{tabular}[t]{l@{}}
	        \multicolumn{1}{@{}l@{}}{\usekomavar*{frombank}:}\\
	        \usekomavar{frombank}
	\end{tabular}
 }}

\begin{document}


%===================================
% receiver
\begin{letter}{
        Dr. Eva Fischbach \\
	Alstertor 21 \\
        20095 Hamburg
}

\opening{Sehr geehrte Frau Fischbach,}

Nach eingehender Pr\"ufung und R\"ucksprache mit dem Rest des
Vorstandes von GNUnet e.V. muss ich Ihnen mitteilen dass wir Ihre
Abmahnung entschieden zur\"uckweisen, da GNUnet e.V. nicht am Hosting
der Webseite \texttt{taler.net} beteiligt ist und auch keinen
Einfluss auf die Registrierung der Bildmarke beim Marken- und Patentamt
(EUTM 17871948) hat.  Weiterhin hat der Verein auch nicht die Absicht,
die Bildmarke von Taler zu verwenden.

Wir verstehen jedoch Ihre Verwirrung, da GNUnet e.V., als eine
Organisation die an der Softwareentwicklung f\"ur Taler mitgewirkt hat,
auf der Seite erw\"ahnt wird. Daher w\"urden wir Ihnen
entgegenkommen und keine Kosten f\"ur unberechtigte Abmahnung
in Rechnung stellen, sofern Sie ihre Abmahnung bis zum 7.5.2018 18 Uhr
zur\"uckziehen.  Ansonsten m\"ussten wir unsererseits einen Anwalt
beauftragen, und Ihnen ggf. die uns daraus entstehende Kosten in
Rechnung stellen.

Das \texttt{taler.net} bei Inria gehostet ist, sehen Sie
wie folgt:
\begin{verbatim}
$ nslookup taler.net
Name:	taler.net
Address: 131.254.145.3

$ nslookup 131.254.145.3
3.145.254.131.in-addr.arpa
name = tripwire.decentralise.rennes.inria.fr.
\end{verbatim}

Wie bereits in unserem Telefongespr\"ach erw\"ahnt, gibt es auch noch
\texttt{gnutaler.org}.  Wir versichern hiermit explizit dass wir diese
Domaine nicht betreiben, und dass uns die Identit\"at der Betreiber
oder des Betreibers unbekannt ist.  Unter \texttt{gnutaler.org} findet
ein Betrugsversuch unter Verletzung von Pers\"onlichkeitsrechten von
GNUnet e.V. Mitgliedern sowie auch von Copyrights von GNUnet e.V.,
Taler Systems SA und Inria statt.  Als Verein sind wir jedoch nicht in
der Lage diese Verletzung effektiv zu verfolgen.

Ein Vorgehen der DTAG gegen die tats\"achlichen Betreiber von
\texttt{gnutaler.org} w\"urde daher von GNUnet e.V. explizit
begr\"usst.  Da die DTAG der Meinung ist, dass das Taler-Logo auf
\texttt{gnutaler.org} dem DTAG-Logo verwechselt werden k\"onnte,
w\"are aus unserer Perspektive ein Vorgehen gegen diese kriminelle
Seite sicherlich auch im Interesse ihrere Mandantin.


Mit freundichen Gr\"ussen

\vspace{3cm}
Christian Grothoff\\
President, GNUnet e. V.


\end{letter}
\end{document}
